\documentclass[11pt]{article}
\usepackage[T1]{fontenc}
\usepackage{lmodern,microtype,amsmath,amssymb,amsthm,booktabs}
\usepackage[margin=1in]{geometry}
\usepackage[hidelinks]{hyperref}
\newtheorem{theorem}{Theorem}[section]
\newtheorem{lemma}[theorem]{Lemma}
\newcommand{\Z}{\mathbb Z}\newcommand{\N}{\mathbb N}
\newcommand{\xor}{\mathbin{\triangle}}
\title{One Bit Off the Highway:\\Uniform Recovery for Langton's Ant}
\author{Hao Ke}
\date{16 October 2026 --- Review draft}
\begin{document}\maketitle
\begin{abstract}
We study a single colour perturbation applied at an arbitrary time and site
along a canonical Langton highway, retaining its entire accumulated wake.
An all-age coupling reduction expresses recovery through 104 guarded
reference phases and finitely many younger exceptions. The resulting
certificate specification has 739 cases, with a reported uniform bound of
64,584 updates and radius 78 after first impact. The same reduction yields
an exact first-impact query and a compressed complete endpoint, with
polynomial bit complexity in the time and perturbation coordinates.
This manuscript develops the analytic reduction and specifies the finite
certificate needed to establish the numerical recovery claims.
\end{abstract}

\section{A growing highway and a single perturbation}
The ant turns right on white and left on black, toggles the read cell, and
moves one unit. Let $F$ be this reversible update on a complete finite black
field, position and heading. Start from the thirteen-black seed
\[
\begin{split}
 P=\{&(0,-2),(0,-1),(0,0),(1,-1),(2,-1),(2,1),(3,-2),\\
     &(3,1),(4,-2),(4,0),(5,-1),(5,0),(6,0)\},
 \qquad X_0=(P,(0,0),\mathrm S).
\end{split}
\]
Its standard word $W$ and drift $d=(2,-2)$ give $X_t=F^tX_0=(B_t,p_t,h_t)$.
We perturb one bit at an arbitrary site $q$:
\[
 Y(t,q)=(B_t\xor\{q\},p_t,h_t),\qquad t\ge0,\ q\in\Z^2.
\]
Thus an existing black cell can be removed as well as a white cell added.
All other cells of the growing canonical history remain present.

The mathematical question is uniform recovery after the first physical
encounter with the perturbation. Its waiting time can depend on an arbitrarily
distant $q$. The reduction below separates this exact waiting-time query from
the finite scattering problem and the all-age extension guard.

The following certificate target records the numerical statement of the
research proof. The analytic implications are proved here; the release
record tracks the corresponding finite data separately.

\begin{theorem}[Certificate implication]\label{thm:recovery}
Assume the finite records and guards specified in Section~\ref{sec:certificate}
pass their complete checks. Then every $Y(t,q)$ eventually follows standard
P104. If $s\ge t$ is the first update at which the unperturbed orbit reads
$q$, recovery is certified by time $s+64584$. From impact to that boundary
the perturbed head stays within Chebyshev distance 78 of $q$. If there is
no such $s$, its entire future read trace agrees with the unperturbed one.

Current-cell perturbations have at most 33,586 future intrinsic $J_{10}$
events, attained at $t=800$. Arbitrary-site perturbations have at most 33,595.
At each fixed phase, all canonical ages of at least 25 periods have the same
normalized infinite response to a current-cell flip.
\end{theorem}
An intrinsic $J_{10}$ event tests the actual nine inverse turns and the current
turn against the cyclic ten-factors of $W$. This definition retains the history
across a perturbation, which may differ from the original canonical past.

\section{The complete finite wake at arbitrary age}
Let $q_i$ be the pre-update positions of the 104 phases and $T_r$ the XOR set
of the first $r$ read positions. The period write set is $T=T_{104}$. The
finite period identity has deposited set
\[
\begin{split}
 D=\{&(0,-3),(0,0),(1,-4),(1,-3),(1,-2),(2,-1),\\
     &(2,0),(3,0),(6,-1),(7,-1),(7,0),(8,-1)\}.
\end{split}
\]
Specifically, $F^{104}X_0$ has field $(P+d)\mathbin{\dot\cup}D$ and pose
$(d,\mathrm S)$, equivalently $D=P\xor(P+d)\xor T$. Distinct points of $D$
lie in distinct full classes modulo $\Z d$, and $D-kd$ is disjoint from $P$
for every $k\ge1$. These are finite word and integer-ray checks.

Translate the state at time $104n$ by $-nd$. Its complete field is
\begin{equation}\label{eq:source}
 C_n=P\mathbin{\dot\cup}\bigcup_{k=1}^n(D-kd),\qquad |C_n|=13+12n.
\end{equation}
Telescoping the period write identity proves this formula by induction.
At phase $r$ the field is $C_n\xor T_r$ and the pose is $(q_r,h_r)$.
For $n\ge N$, two such sources, with or without their current cell toggled,
differ exactly on the finite set $\bigcup_{k=N+1}^n(D-kd)$.

\section{Extending one recovery to every older source}
\begin{lemma}[First-difference coupling]\label{lem:coupling}
If two complete fields at the same pose differ on $E$, and one trajectory
never reads $E$, their entire read and pose traces agree; their field
difference remains $E$ throughout.
\end{lemma}
\begin{proof}
Before a first differing read, both current colours agree. They turn together
and toggle the same site, preserving the difference. Induction rules out such
a first read.
\end{proof}
For an age $N$ define the exclusion set $A_N=\bigcup_{k\ge N+1}(D-kd)$.
It is a guard region; every physical source in~\eqref{eq:source} is finite.

\begin{lemma}[Guarded age extension]\label{lem:age}
Suppose the current-cell perturbation at age $N$, phase $r$, reaches a
certified standard suffix. If its finite recovery footprint and all reads
of that suffix avoid $A_N$, the same normalized infinite trace occurs at
every age $n\ge N$. If its nine inverse reads also avoid $A_N$, all intrinsic
event counts agree.
\end{lemma}
\begin{proof}
The actual finite difference of the sources is contained in $A_N$.
Apply Lemma~\ref{lem:coupling} forwards and, for the history assertion,
to the reversible rule for nine steps backwards.
\end{proof}
The infinite guard has a finite exact test. For a finite read $p$, solve
$e-p=kd$ for each $e\in D$, with $k\ge N+1$. For a suffix phase ray
$v_i+je_{\rm out}$, solve
\begin{equation}\label{eq:guard}
 v_i+je_{\rm out}=e-kd,\qquad j\ge0,\quad k\ge N+1.
\end{equation}
Nonparallel drifts give a unique rational solution by determinants, checked
for integrality and the two inequalities. For equal parallel drifts and
$v_i-e=ad$, a solution exists exactly when $-a\ge N+1$. For opposite drifts,
equal full classes yield an eventual intersection. Distinct full classes
never intersect. Thus every infinite ray and every old-wake epoch is covered.

\section{The finite certificate specification}\label{sec:certificate}
For each phase $r$, use the following age threshold $N_r$, in phase order:
\begin{quote}\small\ttfamily
8,7,9,2,1,5,1,3,5,15,7,2,2,1,2,6,1,2,5,2,3,11,21,3,4,3,\par
2,1,13,17,1,17,3,7,1,3,18,10,7,0,14,8,3,5,9,9,4,4,9,4,\par
5,25,9,13,2,12,7,2,0,16,9,14,12,14,0,3,2,4,14,3,4,8,11,\par
6,12,5,12,4,4,15,14,0,0,15,1,15,0,8,0,1,0,3,1,4,2,5,8,\par
1,3,0,3,2,1,1.
\end{quote}
The certificate must contain exactly the pairs $(n,r)$ with $0\le n\le N_r$:
635 younger exceptions and 104 guarded bases. Each record specifies recovery
length, endpoint pose, outgoing standard phase, exact intrinsic defect count
and last defect time. Its checker must perform these steps:
\begin{enumerate}
\item Reconstruct the complete source~\eqref{eq:source}, its phase writes and
the current-cell flip; recover the actual nine inverse turns.
\item Replay every recovery update and compare the complete endpoint, pose,
read word, event count and maximum distance from the impact cell.
\item Verify the entire proposed standard future by the parity test below,
including its nine-letter incoming history.
\item At $n=N_r$, verify the finite, inverse and infinite footprints against
$A_{N_r}$ using~\eqref{eq:guard}.
\end{enumerate}

At suffix phase $i$ and epoch $k$, the necessary actual colour is
\begin{equation}\label{eq:tail}
 1_B(v_i+ke_{\rm out})\ \mathbin{\oplus}\
 \#\{(j,z):v_j-v_i=ze_{\rm out},\ 0\le z\le k,
                 \ z>0\text{ or }(z=0,j<i)\}\pmod2.
\end{equation}
Here $B$ is the entire endpoint field. A black cell hits a phase ray at most
once, while each same-ray earlier phase supplies one threshold. Epoch zero,
every hit and its successor, every self-write threshold, and the final
infinite interval therefore form a finite complete check of~\eqref{eq:tail}.
Induction over suffix time proves permanence.

The recorded aggregate targets are maximum recovery length 64,584 at
$(n,r)=(22,51)$, maximum radius 78 at $(25,51)$, and maximum event count
33,586 at $(7,72)$. The last corresponds to $t=104\cdot7+72=800$.
These are bounds at the certificate's recovery boundaries; the time bound
is not a claim about the latest first permanent entry.

\begin{proof}[Current-cell part of Theorem~\ref{thm:recovery}]
The explicit records cover $n<N_r$. Each guarded base covers all $n\ge N_r$
by Lemma~\ref{lem:age}. Taking the record maxima proves the bounds and the
attained event count. Since $N_r\le25$, responses stabilize by age 25 at
every phase.
\end{proof}

\section{Arbitrary perturbation sites and exact endpoints}
For arbitrary $t,q$, its first future canonical read is the minimum of
\[
 \{104n+i\ge t:q=q_i+nd,\ n\ge0,\ 0\le i<104\}.
\]
Each phase supplies at most one $n$. Thus 104 exact collinearity and
divisibility checks decide whether impact occurs and give its exact time.
If there is none, coupling proves agreement forever. Otherwise the two
orbits agree up to impact $s$, when the perturbed state is precisely $X_s$
with its current cell toggled. The current-cell result now applies.
Before impact all forward letters are canonical; only the first nine tests
can see altered inverse history. Adding those nine to the post-impact bound
gives 33,595 and completes Theorem~\ref{thm:recovery}.

Write $s=104n+r$, choose $m=\min(n,N_r)$, and replay the bounded reference
record at $(m,r)$. If its normalized write set is $Z$ and endpoint pose is
$(p_{\rm end},h_{\rm end})$, the true complete endpoint is
\[
 B_s\xor\{q\}\xor(Z+nd),\qquad
 (p_{\rm end}+nd,h_{\rm end}).
\]
Formula~\eqref{eq:source} represents $B_s$ by twelve finite arithmetic
segments, the fixed seed and phase writes. Its binary size is polynomial in
the input time and coordinates. The number of reference updates is bounded
by an absolute constant. This proves the exact polynomial-bit-time endpoint
algorithm, conditional on the same finite certificate.

\section{Relation to existing results}
The one-black theorem classifies complete inputs with one initially black
cell. The present family instead contains $13+12n$ canonical historical
cells before its phase writes, with one arbitrary colour toggle. The
all-age guard is what extends a finite scattering check to those unbounded
histories. The factor paper supplies the intrinsic event language; its
local predicate is used here to count the recovery rather than to predict
all possible initial fields.

A recovered standard trace can leave a finite scattering scar. Repeated
perturbations therefore require an additional closure theorem for the
complete repaired fields. The single-toggle family developed here supplies
a concrete interface for that further problem.

The stable-read and symbolic-word viewpoints are prior methods
\cite{gajardo,jillhewar}. This manuscript's specific content is the full
canonical-age quantifier, guarded recovery reduction and finite certificate
specification. The repository review record identifies which data have been
replayed and which must still be restored before numerical certification.

\appendix
\section{Canonical word}
The following concatenated lines specify the 104-letter word:
\begin{quote}\ttfamily\small
LRLLRRRRLLRLLRRRRLLRRRRLLRLRRRRLRLLLLRRRRLRRLRRRRLLLL\par
RLRRRRLRRRRLLLLRLRRRRLRLLRRLLLLRRLLRRRRLLRRLRLLRLLR.
\end{quote}
The period and deposition identities can be checked from these letters and
$P$ directly. The required finite certificate and source manifest are
identified in the accompanying evidence record.

\begin{thebibliography}{9}
\bibitem{gajardo} A.~Gajardo. A symbolic projection of Langton's Ant.
\emph{Discrete Mathematics \& Theoretical Computer Science}, 2003.
\href{https://doi.org/10.46298/dmtcs.2312}{doi:10.46298/dmtcs.2312}.
\bibitem{jillhewar} A.~Jillhewar. \emph{Finite-Support Periodic Highways of
Langton's Ant: Necessary Conditions, Transverse Exclusions, and Exact Search},
2026. Repository revision \texttt{76a7e234e577d992d086c204c3dea00ae72cf10a},
\url{https://github.com/Atharva12456/langtons-ant-highway}.
\end{thebibliography}
\end{document}
